% ═══════════════════════════════════════════════════════════════
%  PART XVII — APRIL 8, 2026: GAUSSIAN UNIFICATION + SELECTION
% ═══════════════════════════════════════════════════════════════
\clearpage
\part{\texorpdfstring{The Gaussian Unification: All Couplings from $z=11+4i$}{The Gaussian Unification: All Couplings from z=11+4i}}

\section{\texorpdfstring{Lock~0: The Structural Selection of $q=3$}{Lock 0: The Structural Selection of q=3}}
\label{sec:lock0}

The selection arguments in this Part are \emph{fits}: they take a measured constant
and ask which $q$ reproduces it.  They are stated as such.  There is also a
selection argument that uses no measured quantity at all, and it is placed first
because it is of a different kind.

\begin{theorem}[The substrate tower has exactly one sibling]
\label{thm:springer-tower}
Let $w\in W(E_8)$ be a regular element of order $d$ in Springer's sense.  Then
$C_{W(E_8)}(w)$ is the complex reflection group whose degrees are the degrees of
$W(E_8)$ divisible by $d$, so $|C| = \prod\{d_i : d\mid d_i\}$.  Over the degrees
$2,8,12,14,18,20,24,30$ this gives a whole tower, and \emph{only three values of
$d$ yield a derived subgroup transitive on the $240$ roots}:
\[
\begin{array}{c|c|c|c}
d & |C| & \text{blocks} & \text{base}\\\hline
3,\,6 & 155\,520 = 12\cdot18\cdot24\cdot30 & [2,3,6] & 40 = W(3,3)\\
4 & 46\,080 = 8\cdot12\cdot20\cdot24 & [2,4,16] & 15 = W(2,2)
\end{array}
\]
with $d=3$ and $d=6$ the same tower.  So the Eisenstein construction that produces
$W(3,3)$ is \textbf{not unique}: there is exactly one sibling, the Gaussian
$d=4$ tower ($\mathrm{ST}_{31}$), and its base is the doily.
\end{theorem}

\noindent
This forbids one argument that might have been tempting --- $q=3$ cannot be
selected merely by observing that the Eisenstein tower exists.  It also makes the
selection question precise, because the two towers are now explicitly comparable.
They differ in exactly two ways, and neither is a fitted quantity.

\begin{theorem}[Contextual selection]
\label{thm:contextual-selection}
An ovoid of $\GQ(q,q)$ is precisely a Kochen--Specker $0/1$ colouring of its
$st{+}1$ contexts, and by Thas' theorem $W(q,q)$ has ovoids if and only if $q$ is
even.  Counted by exact cover:
\[
W(2,2):\ 6\ \text{ovoids},\qquad W(3,3):\ 0\ \text{ovoids}.
\]
Hence the Gaussian base admits a global noncontextual value assignment and the
Eisenstein base does not.  \emph{Of the two towers the construction actually
produces, only $q=3$ is contextual.}
\end{theorem}

\begin{theorem}[Shape of the section obstruction]
\label{thm:obstruction-shape}
Both towers are obstructed: neither $240\to$ base admits an equivariant section.
They are obstructed differently.  The Eisenstein fibre is
$\langle c^5\rangle = \langle -1,\omega\rangle \cong C_6 = C_2\times C_3$, whose
obstruction splits into two \emph{independent} coordinates --- a sign and a phase,
separated by the centre of $\Sp(4,3)$.  The Gaussian fibre is $C_4$, which is
cyclic with a unique subgroup of order two and therefore admits no complement: its
obstruction is one indecomposable class.  Only the Eisenstein tower carries two
independent obstruction coordinates.
\end{theorem}

\begin{remark}
Theorems~\ref{thm:contextual-selection} and~\ref{thm:obstruction-shape} use no
measured constant.  They select $q=3$ from a two-element list by structure alone,
which the Locks below do not do.  The Locks remain what they are --- consistency
checks against experiment, and evidence that the arithmetic is not accidental ---
but the structural statement is logically prior and should be read first.
Witnesses: \texttt{analysis/w33\_pass1039\_springer\_tower.g},
\texttt{...\_pass1039b\_gaussian\_base.g},
\texttt{...\_pass1042\_tower\_discriminators.g}.
\end{remark}


\section{\texorpdfstring{Lock~9: $\alpha^{-1}=137$ Uniquely Selects $q=3$}{Lock 9: alpha-inverse equals 137 uniquely selects q=3}}

For $\GQ(q,q)$, $(k{-}1)^2+\mu^2 = q^4+2q^3+2$.
Setting this equal to~137:
\begin{equation}
\boxed{q^3(q+2) = 135 = 3^3\times 5.}
\end{equation}
This has the \textbf{unique} positive integer solution $q=3$.
Verified computationally for $q\in\{2,3,4,5,7,8,9,11,13\}$:
no other $q$ gives~$137$.

\begin{theorem}[7/7 Observable Test]
Among all $\GQ(q,q)$, only $q=3$ simultaneously matches:
$\alpha^{-1}\approx 137$,
$\sin^2\theta_{12}\approx 0.307$,
$\sin^2\theta_{23}\approx 0.546$,
$\sin^2\theta_{13}\approx 0.022$,
$\alpha_s\approx 0.118$,
$f = 24$,
$E = 240$.
No other $q$ matches even one observable.
\end{theorem}


\section{\texorpdfstring{Lock~10: The Gaussian Coincidence $k{-}1=\Phi_3{-}\lambda$}{Lock 10: The Gaussian Coincidence k-1 = Phi3-lambda}}

For general $\GQ(q,q)$: $k{-}1=q^2{+}q{-}1$ and $\Phi_3{-}\lambda=q^2{+}2$.
These are equal if and only if $q=3$.

\begin{theorem}[Gaussian Unification]
\label{thm:gaussian}
At $q=3$, the Gaussian integer $z = (k{-}1)+\mu i = (\Phi_3{-}\lambda)+\mu i
= 11+4i$ encodes all three gauge couplings:
\begin{align}
\alpha_{\mathrm{em}}^{-1} &= |z|^2 = 137 \quad (\text{tree level}), \\
\alpha_s &= \mu^2/|z|^2 = 16/137 \quad (\text{tree level, $1.3\sigma$}), \\
\sin^2\theta_W &= q/\Phi_3 = q/((k{-}1)+\lambda) \quad (\text{dressed}).
\end{align}
The coupling ratios at tree level satisfy $\alpha_s/\alpha_{\mathrm{em}} = \mu^2 = s^2$
and
\begin{equation}
\alpha^{-1} = \left(\frac{q}{\sin^2\theta_W} - \lambda\right)^2 + \mu^2.
\end{equation}
This algebraic relation between $\alpha^{-1}$ and $\sin^2\theta_W$
holds \textbf{only at $q=3$}.
\end{theorem}


\section{The 7/7 Definitive Selection Table}

\begin{center}
\begin{tabular}{@{}crrrrrrrr@{}}
\toprule
$q$ & $\alpha^{-1}_0$ & $\sin^2\theta_{12}$ & $\sin^2\theta_{23}$
& $\sin^2\theta_{13}$ & $\alpha_s$ & $f$ & $E$ & score \\
\midrule
$2$ & $34$ & $0.429$ & $0.429$ & $0.059$ & $0.204$ & $9$ & $45$ & $0/7$ \\
$\mathbf{3}$ & $\mathbf{137}$ & $\mathbf{0.308}$ & $\mathbf{0.538}$
& $\mathbf{0.022}$ & $\mathbf{0.118}$ & $\mathbf{24}$ & $\mathbf{240}$
& $\mathbf{7/7}$ \\
$4$ & $386$ & $0.238$ & $0.619$ & $0.009$ & $0.077$ & $50$ & $850$ & $0/7$ \\
$5$ & $877$ & $0.194$ & $0.677$ & $0.004$ & $0.054$ & $90$ & $2340$ & $0/7$ \\
$7$ & $3089$ & $0.140$ & $0.754$ & $0.000$ & $0.031$ & $224$ & $11200$ & $0/7$ \\
\bottomrule
\end{tabular}
\end{center}
\emph{Seven independent observables from seven experiments.
Only $q=3$ matches all seven. No other $q$ matches even one.}
